% SPDX-License-Identifier: GPL-3.0-or-later
% Copyright (C) 2026 Alessandro Greco
\chapter{От луковой маршрутизации NRL к Tor Project}\label{cap:03-nascita-tor}
\aggiornatoal{23 сентября 2026 г.}{Syverson, \emph{A Peel of Onion} (2011); страница истории Tor Project}

\begin{obiettivi}
К концу этой главы вы сможете:
\begin{itemize}[nosep]
  \item проследить происхождение onion routing в Naval Research Laboratory и описать три поколения этой технологии;
  \item объяснить, чем Tor отличается от предыдущих поколений;
  \item перечислить основные этапы истории Tor Project с 2002 года до наших дней;
  \item объяснить происхождение названия \enquote{Tor}.
\end{itemize}
\end{obiettivi}

% ---------------------------------------------------------------------------
\section{В Naval Research Laboratory}\label{sec:nascita-nrl}

В 1995 году трое исследователей Naval Research Laboratory (NRL) ВМС США --- Дэвид Голдшлаг, Майкл Рид и Пол Сайверсон --- задались вопросом, можно ли устанавливать интернет-соединения, не раскрывая, кто с кем общается~\cite{torhistory}. Их целью была не анонимность сама по себе, а, как мы уже видели, стремление \enquote{отделить идентификацию от маршрутизации} (\cref{sec:anon-termini}). Первый вариант архитектуры они представили на \emph{Information Hiding Workshop} в мае 1996 года~\cite{goldschlag1996hiding}. До конца того же года появился второй вариант, в котором пользовательское приложение, строящее цепи, было отделено от узлов, передающих трафик~\cite{syverson2011peel}. Уже в 1996 году лаборатория получила разрешение опубликовать исходный код, ещё до широкого распространения самой идеи \emph{программного обеспечения с открытым исходным кодом}. Разработчики понимали: никто не станет доверять сети, код которой нельзя изучить.

Название происходит от структуры данных: сообщение обёрнуто слоями шифрования, каждый из которых снимается очередным узлом, словно слои луковицы.

% ---------------------------------------------------------------------------
\section{Три поколения}\label{sec:nascita-generazioni}

Сайверсон различает три поколения разработок NRL~\cite{syverson2011peel}.

\begin{figure}[htbp]
\centering
% SPDX-License-Identifier: GPL-3.0-or-later
% Copyright (C) 2026 Alessandro Greco
%% Русское альтернативное описание EPUB.
\altfig{Сравнение трёх поколений луковой маршрутизации. В первых двух клиент заранее создаёт одну многослойно зашифрованную структуру для всего маршрута; в Tor цепь расширяется по одному переходу, с отдельным обменом ключами для каждого ретранслятора.}
% Три поколения луковой маршрутизации NRL.
\begin{tikzpicture}[font=\sffamily\scriptsize,
  nodo/.style={relay,minimum size=0.65cm,font=\sffamily\tiny},
  eti/.style={anchor=east,align=right,font=\sffamily\small}]
  \node[eti] at (-0.5,0) {\textbf{1-е покол.}\\(1996)};
  \node[client,minimum width=1.1cm,minimum height=0.6cm,font=\sffamily\tiny,align=center] (c1) at (0.6,0) {клиент\\+ узел};
  \foreach \i in {1,...,5} \node[nodo] (a\i) at (0.6+1.05*\i,0) {};
  \draw[cellaf] (c1) -- (a1); \foreach \i/\j in {1/2,2/3,3/4,4/5} \draw[cellaf] (a\i) -- (a\j);
  \draw[decorate,decoration={brace,amplitude=3pt}] (1.3,0.45) -- node[above=3pt,font=\sffamily\tiny,align=center] {5 фиксированных узлов; одна многослойная структура\\задаёт всю цепь} (5.95,0.45);
  \node[eti] at (-0.5,-1.6) {\textbf{2-е покол.}};
  \node[client,minimum width=1.1cm,minimum height=0.6cm,font=\sffamily\tiny] (c2) at (0.6,-1.6) {клиент};
  \foreach \i in {1,...,4} \node[nodo] (b\i) at (0.6+1.05*\i,-1.6) {};
  \draw[cellaf] (c2) -- (b1); \foreach \i/\j in {1/2,2/3,3/4} \draw[cellaf] (b\i) -- (b\j);
  \node[anchor=west,font=\sffamily\tiny,align=left] at (5.4,-1.6) {клиент отделён от узла,\\переменная длина, перемешивание (отказались)};
  \node[eti] at (-0.5,-3.2) {\textbf{3-е покол.}\\Tor (2002)};
  \node[client,minimum width=1.1cm,minimum height=0.6cm,font=\sffamily\tiny] (c3) at (0.6,-3.2) {клиент};
  \foreach \i in {1,2,3} \node[nodo] (d\i) at (0.6+1.05*\i,-3.2) {};
  \draw[cellaf] (c3) -- node[above,font=\sffamily\tiny] {1} (d1);
  \draw[cellaf] (d1) -- node[above,font=\sffamily\tiny] {2} (d2);
  \draw[cellaf] (d2) -- node[above,font=\sffamily\tiny] {3} (d3);
  \node[anchor=west,font=\sffamily\tiny,align=left] at (4.3,-3.2) {3 узла, построение \emph{по одному переходу}:\\отдельный обмен ключами, прямая секретность};
\end{tikzpicture}

\caption{Три поколения onion routing. Ключевое изменение в Tor состояло не в числе переходов, а в способе построения цепи: вместо подготовки клиентом одной полной «луковицы» цепь расширяется по одному переходу, с отдельным обменом ключами для каждого узла.}
\label{fig:generazioni}
\end{figure}

\paragraph{Первое поколение.} Использовались цепи фиксированной длины из пяти узлов, предназначенные, в частности, для соединения защищённых сетей двух организаций. При трёх переходах враждебный промежуточный узел сразу видел бы обе взаимодействующие сети. При четырёх атакующий, скомпрометировав узел рядом с одной конечной точкой, знал бы, какой ещё узел нужно атаковать, чтобы восстановить всю цепь. При пяти такой короткий путь был невозможен без предварительной компрометации ещё одного узла. Клиент был встроен в сам узел. Сеть была статической: двадцать или сто узлов под управлением постоянных организаций, конфигурации и ключи распространялись вручную. Каждому приложению требовался собственный прокси --- для веб-сёрфинга, удалённого входа или электронной почты, --- поскольку SOCKS тогда ещё не был широко распространён. Цепь строилась как единая «луковица», зашифрованная открытыми ключами всех узлов.

\paragraph{Второе поколение.} Клиент отделили от узла, длина цепи стала переменной, были введены механизмы перемешивания и заполнения между цепями, от которых впоследствии отказались. Распространение информации о сети было предусмотрено, но полностью так и не реализовано.

\paragraph{Третье поколение: Tor.} Самым важным изменением стало построение цепи \textbf{по одному переходу}, с обменом Диффи--Хеллмана с каждым узлом (\cref{sec:circ-telescopio}) --- идеей, заимствованной у проекта Cebolla. Это обеспечивает прямую секретность и позволяет проверять каждый переход по мере расширения цепи. Tor также добавил серверы каталогов, поддержку SOCKS, скрытые сервисы с точками рандеву и, прежде всего, архитектуру, рассчитанную на развёртывание и эксплуатацию добровольцами по всему миру.

\begin{storia}[Название]
Заголовок статьи 2004 года \titolo{Tor: The Second-Generation Onion Router} породил два заблуждения, которые позднее признал сам Сайверсон~\cite{syverson2011peel}. Tor относится к \emph{третьему}, а не ко второму поколению, и буква \enquote{r} в Tor не означает \emph{router}. Когда Роджер Динглдайн начал работать над проектом NRL, существовало множество систем onion routing. На вопрос \enquote{какая именно?} он отвечал \enquote{\emph{the onion routing}}, имея в виду исходную систему. Рэйчел Гринстадт заметила, что первые буквы образуют удачный акроним. Динглдайн также обратил внимание на рекурсивное прочтение \emph{Tor's onion routing} и решил, что название следует писать \enquote{Tor}, а не \enquote{TOR}: это обычное слово, которое по-немецки также означает \emph{ворота}, например городские. Во всей остальной книге мы пишем \emph{Tor}, когда речь идёт о сети и проекте, и \emph{tor} со строчной буквы, когда речь идёт о программе.
\end{storia}

% ---------------------------------------------------------------------------
\section{Сеть Tor}\label{sec:nascita-rete}

Роджер Динглдайн присоединился к проекту NRL в начале 2000-х годов, а вскоре к нему присоединился Ник Мэтьюсон. Сеть Tor начала работу в \textbf{октябре 2002 года}, причём исходный код с самого начала распространялся под лицензией свободного программного обеспечения. К концу 2003 года в сети было около дюжины узлов, управляемых добровольцами, почти все в США и один в Германии~\cite{torhistory}. В августе 2004 года Динглдайн, Мэтьюсон и Сайверсон представили проект на USENIX Security Symposium~\cite{dingledine2004tor}.

Основные последующие этапы были следующими~\cite{torhistory}:
\begin{center}\small
\begin{tabularx}{\textwidth}{@{}l>{\raggedright\arraybackslash}X@{}}
\toprule
Год & Событие \\
\midrule
2002 & Сеть Tor начинает работу; исходный код свободно доступен \\
2004 & Electronic Frontier Foundation начинает финансировать работу Динглдайна и Мэтьюсона \\
2006 & Tor Project, Inc. учреждается как некоммерческая организация США (501(c)(3)) \\
2007 & Начинается разработка устойчивых к цензуре мостов (\cref{cap:16-censura}) \\
2008 & Начинается разработка Tor Browser (\cref{cap:23-browser}) \\
2010--2011 & Tor становится важным инструментом во время восстаний «арабской весны» \\
2013 & Публикации Эдварда Сноудена о массовой слежке, в которых Tor играет центральную роль \\
2017 & Onion-сервисы v3 (\cref{cap:14-onion-services}) \\
2021 & Onion-сервисы v2 отключены \\
2022--2023 & Управление перегрузкой (\cref{sec:cc}) и proof of work для onion-сервисов (\cref{cap:15-dos}) \\
2024 & WebTunnel; перенос распределения мостов на rdsys \\
2025--2026 & CGO в tor (\cref{sec:cgo}); Arti как клиент и для onion-сервисов (\cref{cap:30-arti}) \\
\bottomrule
\end{tabularx}
\end{center}

% ---------------------------------------------------------------------------
\section{Организация и сообщество}\label{sec:nascita-comunita}

Tor Project, Inc.\ разрабатывает программное обеспечение, поддерживает спецификации, совместно с независимыми операторами обслуживает доверенные серверы каталогов (directory authorities), измеряет сеть и поддерживает сообщество. Но организация не владеет сетью: тысячи добровольцев --- частные лица, университеты и ассоциации --- эксплуатируют реле в десятках стран. Это архитектурное решение учитывает опыт Freedom (\cref{sec:radici-bassa}): ни одна организация не должна иметь возможность единолично контролировать сеть или быть вынужденной скомпрометировать её пользователей.

Финансирование организации неоднократно становилось предметом публичных споров. На протяжении многих лет средства поступали от государственных ведомств, особенно в США, частных фондов и отдельных жертвователей. Организация ежегодно публикует финансовую отчётность. Вопрос \enquote{Можно ли доверять программному обеспечению, которое частично финансировалось государством?} прежде всего технический, а не политический: код и спецификации публичны и проверяемы, архитектурные решения открыто обсуждаются в предложениях Tor (proposal) (\cref{cap:31-futuro}), а безопасность системы не зависит от доверия к Tor Project. Вместо этого она опирается на распределение сети между множеством независимых операторов (\cref{cap:04-modello-minaccia}).

\begin{approfondimento}[Публикации 2013 года]
Среди документов, опубликованных Сноуденом в 2013 году, были внутренние презентации NSA о Tor. В них признавалась сложность систематической деанонимизации пользователей Tor и описывались целевые атаки на отдельных пользователей, эксплуатировавшие уязвимости браузера, а не самой сети. Это косвенно подтверждает принцип, к которому мы будем неоднократно возвращаться: самые опасные слабости нередко находятся не в самом протоколе, а в окружающих его системах (\cref{cap:28-opsec}).
\end{approfondimento}

% ---------------------------------------------------------------------------
\section*{Итоги}
\begin{itemize}
  \item Onion routing появился в 1995 году в Naval Research Laboratory благодаря Голдшлагу, Риду и Сайверсону с целью отделить идентификацию от маршрутизации.
  \item Первое поколение использовало фиксированные цепи из пяти узлов и статическую сеть. Tor, третье поколение, строит цепи по одному переходу, обеспечивает прямую секретность и использует каталоги и точки рандеву.
  \item Название \enquote{Tor} происходит от \emph{the onion routing}; заглавной является только первая буква.
  \item Сеть работает с 2002 года; Tor Project, Inc.\ существует с 2006 года; реле эксплуатируются независимыми добровольцами.
\end{itemize}

% ---------------------------------------------------------------------------
\section*{Упражнения}

\begin{esercizio}[1]\label{es:03-cinque}
Почему первое поколение использовало цепи из пяти узлов? Восстановите приведённые в этой главе аргументы для трёх и четырёх переходов и объясните, почему Tor использует только три.
\end{esercizio}

\begin{esercizio}[2]\label{es:03-cipolla}
Сравните построение цепи как единой «луковицы» (первое поколение) с телескопическим построением Tor с точки зрения прямой секретности и возможности немедленно определить не отвечающий узел.
\end{esercizio}

\begin{esercizio}[2]\label{es:03-fiducia}
Напишите краткий аргумент объёмом около одной страницы в ответ на вопрос \enquote{Можно ли доверять Tor с учётом того, кто его финансировал?}. Разделите то, что можно проверить технически, и то, что действительно требует доверия.
\end{esercizio}

% ---------------------------------------------------------------------------
\section*{Дополнительная литература}
\begin{itemize}
  \item P.~Syverson, \titolo{A Peel of Onion}~\cite{syverson2011peel}.
  \item D.~Goldschlag, M.~Reed, P.~Syverson, \titolo{Hiding Routing Information}~\cite{goldschlag1996hiding}.
  \item Страница истории Tor Project~\cite{torhistory}.
\end{itemize}
